% =====================================================================
% ICBEA 2027 — Reviewer-response LaTeX package
% Generated from revision_*.csv (all CV now runs INSIDE the 80% train split)
% Paste each block into the marked place in the Overleaf source.
% =====================================================================


% ---------------------------------------------------------------------
% [1] REPLACE Table II  (tab:cv) — CV inside training split
% ---------------------------------------------------------------------
\begin{table}[htbp]
\caption{Five-fold stratified cross-validation results (cross-validation performed
within the 80\% training partition). Values are mean $\pm$ standard deviation over folds.}
\label{tab:cv}
\centering
\footnotesize
\begin{tabular}{lcccc}
\toprule
Model & Precision & F1 & Accuracy & Macro-F1 \\
\midrule
XGBoost & \textbf{0.9730}$\pm$0.0030 & \textbf{0.9719}$\pm$0.0033 & \textbf{0.9725}$\pm$0.0032 & \textbf{0.9443}$\pm$0.0184 \\
SVM & 0.9665$\pm$0.0045 & 0.9656$\pm$0.0045 & 0.9663$\pm$0.0044 & 0.9405$\pm$0.0182 \\
Logistic Regression & 0.9625$\pm$0.0059 & 0.9618$\pm$0.0058 & 0.9626$\pm$0.0056 & 0.9290$\pm$0.0206 \\
Random Forest & 0.9542$\pm$0.0042 & 0.9469$\pm$0.0040 & 0.9511$\pm$0.0041 & 0.8879$\pm$0.0066 \\
\bottomrule
\end{tabular}
\end{table}


% ---------------------------------------------------------------------
% [2] REPLACE Table III (tab:kappa) — holdout columns unchanged (they were
%     always leakage-free); CV columns updated
% ---------------------------------------------------------------------
\begin{table}[htbp]
\caption{Holdout and cross-validated Cohen's $\kappa$ and macro-F1
(cross-validation within the 80\% training partition).}
\label{tab:kappa}
\centering
\footnotesize
\begin{tabular}{lcccc}
\toprule
Model & Hold. $\kappa$ & Hold. macro-F1 & CV $\kappa$ & CV macro-F1 \\
\midrule
XGBoost & \textbf{0.9767} & \textbf{0.9534} & \textbf{0.9712}$\pm$0.0033 & \textbf{0.9443}$\pm$0.0184 \\
SVM & 0.9685 & 0.9315 & 0.9647$\pm$0.0046 & 0.9405$\pm$0.0182 \\
Logistic Regression & 0.9641 & 0.9080 & 0.9608$\pm$0.0058 & 0.9290$\pm$0.0206 \\
Random Forest & 0.9546 & 0.8945 & 0.9487$\pm$0.0043 & 0.8879$\pm$0.0066 \\
\bottomrule
\end{tabular}
\end{table}


% ---------------------------------------------------------------------
% [3] NEW — per-class table (reviewer comment 3). Full-width,
%     two side-by-side blocks to save vertical space.
%     Supports are the TRAIN-partition counts (6,651 samples, pooled 5-fold CV).
% ---------------------------------------------------------------------
\begin{table*}[htbp]
\caption{Per-class precision, recall and F1 for XGBoost under pooled 5-fold
cross-validation within the training partition (6{,}651 samples). TCGA project
codes are given for reference.}
\label{tab:perclass}
\centering
\footnotesize
\begin{tabular}{lcccc|cccc}
\toprule
TCGA & N & Prec. & Rec. & F1 & TCGA & N & Prec. & Rec. & F1 \\
\midrule
ACC  & 61  & 0.952 & 0.967 & 0.959 & MESO & 70  & 0.985 & 0.929 & 0.956 \\
BLCA & 321 & 0.975 & 0.981 & 0.978 & OV   & 7   & 1.000 & \textbf{0.429} & 0.600 \\
BRCA & 605 & 0.993 & 0.997 & 0.995 & PAAD & 141 & 0.979 & 0.979 & 0.979 \\
CESC & 233 & 0.978 & 0.949 & 0.963 & PCPG & 129 & 0.992 & 0.977 & 0.984 \\
CHOL & 29  & 0.962 & 0.862 & 0.909 & PRAD & 387 & 0.995 & 1.000 & 0.997 \\
COAD & 228 & 0.869 & 0.956 & 0.910 & READ & 73  & 0.824 & \textbf{0.575} & 0.677 \\
DLBC & 38  & 0.881 & 0.974 & 0.925 & SARC & 200 & 0.933 & 0.980 & 0.956 \\
ESCA & 144 & 0.914 & 0.889 & 0.901 & SKCM & 281 & 0.986 & 0.982 & 0.984 \\
HNSC & 405 & 0.993 & 0.990 & 0.991 & STAD & 292 & 0.943 & 0.956 & 0.949 \\
KICH & 52  & 0.962 & 0.981 & 0.971 & TGCT & 106 & 0.991 & 0.981 & 0.986 \\
KIRC & 245 & 0.971 & 0.955 & 0.963 & THCA & 395 & 1.000 & 0.998 & 0.999 \\
KIRP & 214 & 0.967 & 0.953 & 0.960 & THYM & 95  & 1.000 & 0.979 & 0.989 \\
LAML & 130 & 1.000 & 1.000 & 1.000 & UCEC & 329 & 0.947 & 0.985 & 0.966 \\
LGG  & 405 & 1.000 & 0.998 & 0.999 & UCS  & 44  & 0.907 & 0.886 & 0.897 \\
LIHC & 285 & 0.997 & 0.990 & 0.993 & UVM  & 64  & 0.985 & 1.000 & 0.992 \\
LUAD & 358 & 0.975 & 0.983 & 0.979 & \multicolumn{4}{l}{\textit{Macro avg F1} = 0.9461} \\
LUSC & 285 & 0.972 & 0.961 & 0.967 & \multicolumn{4}{l}{\textit{Weighted avg F1} = 0.9720} \\
\bottomrule
\end{tabular}
\end{table*}


% ---------------------------------------------------------------------
% [4] REPLACE Table IV (tab:wilcoxon) — CORRECTED 2026-11-24.
%     CIs and p-values now both use the Nadeau-Bengio variance
%     (1/J + n_test/n_train) so they are mutually consistent.
%     25 paired folds, J=25, df=24, two-sided alpha=0.05.
% ---------------------------------------------------------------------
\begin{table}[htbp]
\caption{Paired comparison over 25 repeated-CV folds within the training
partition. Mean difference (XGBoost $-$ baseline), 95\% confidence interval,
Cohen's $d$, and Nadeau--Bengio corrected two-sided $p$-value. Intervals and
$p$-values are computed from the same corrected variance.}
\label{tab:wilcoxon}
\centering
\footnotesize
\begin{tabular}{llccccc}
\toprule
Baseline & Metric & Mean $\Delta$ & 95\% CI & $d$ & $p_{\mathrm{NB}}$ \\
\midrule
\multirow{4}{*}{SVM}
 & Weighted precision & +0.0060 & [0.0015, 0.0105] & 1.48 & 0.011 \\
 & Weighted F1        & +0.0060 & [0.0017, 0.0103] & 1.55 & 0.008 \\
 & Accuracy           & +0.0061 & [0.0020, 0.0102] & 1.64 & 0.006 \\
 & Cohen's $\kappa$   & +0.0064 & [0.0021, 0.0107] & 1.64 & 0.006 \\
 & Macro-F1$^{\dagger}$ & +0.0070 & [$-$0.0145, 0.0283] & 0.36 & 0.509 \\
\midrule
\multirow{4}{*}{Logistic Reg.}
 & Weighted precision & +0.0079 & [0.0030, 0.0128] & 1.78 & 0.003 \\
 & Weighted F1        & +0.0077 & [0.0031, 0.0123] & 1.88 & 0.002 \\
 & Accuracy           & +0.0076 & [0.0031, 0.0120] & 1.89 & 0.002 \\
 & Cohen's $\kappa$   & +0.0080 & [0.0033, 0.0126] & 1.90 & 0.002 \\
 & Macro-F1$^{\dagger}$ & +0.0176 & [$-$0.0060, 0.0413] & 0.83 & 0.137 \\
\bottomrule
\multicolumn{6}{p{0.95\columnwidth}}{\footnotesize $^{\dagger}$Not significant after
Nadeau--Bengio correction; the advantage is restricted to weighted metrics.}
\end{tabular}
\end{table}
% NOTE: requires \usepackage{multirow} in the preamble.


% ---------------------------------------------------------------------
% [5] REPLACE Table V (tab:ablation) — leave-one-modality-out WITH paired
%     uncertainty (25 paired folds, Nadeau-Bengio corrected).
% ---------------------------------------------------------------------
\begin{table}[htbp]
\caption{Leave-one-modality-out ablation for XGBoost. Point estimates use the
untouched 20\% holdout and 5-fold CV within the training partition; the last two
columns give the paired difference versus the full model over 25 repeated folds
(Nadeau--Bengio corrected).}
\label{tab:ablation}
\centering
\footnotesize
\begin{tabular}{lcccccc}
\toprule
Configuration & \#Feat. & Hold.\ acc & CV acc & CV m-F1 & $\Delta$m-F1 vs full & $p$ \\
\midrule
Full multi-omics      & 9{,}844 & 0.9778 & \textbf{0.9725} & 0.9443 & --- & --- \\
without mRNA          & 6{,}627 & 0.9790 & 0.9681 & 0.9454 & +0.0064 [$-$0.013, 0.026] & 0.507 \\
without miRNA         & 9{,}461 & 0.9681 & 0.9645 & 0.9259 & +0.0125 [$-$0.004, 0.029] & 0.128 \\
without CNV           & 6{,}739 & 0.9790 & 0.9722 & 0.9479 & $-$0.0016 [$-$0.019, 0.015] & 0.851 \\
without Methylation   & 6{,}705 & 0.9705 & 0.9657 & 0.9296 & +0.0093 [$-$0.001, 0.020] & 0.078 \\
mRNA-only             & 3{,}217 & 0.9603 & 0.9547 & 0.9038 & +0.0368 [0.025, 0.049] & $<$0.001 \\
\bottomrule
\end{tabular}
\end{table}
% Interpretation: only mRNA-only is significantly worse than the full model.
% Removing any single modality does not significantly change macro-F1.


% ---------------------------------------------------------------------
% [6] ADD to end of Sec. III-A (Preprocessing) — leakage statement.
%     This is TRUE for this codebase: CV now runs inside the train split.
% ---------------------------------------------------------------------
To prevent information leakage, the dataset was first partitioned using a
stratified 80/20 train--test split. All data-dependent preprocessing
operations, including missing-value imputation and feature scaling, were fitted
exclusively on the training data. During cross-validation, the scaling
parameters were independently estimated on each training fold and applied to
its corresponding validation fold, and all cross-validation was performed
strictly within the 80\% training partition. The 20\% holdout set was kept
separate from model selection and cross-validation and was used only for the
final evaluation. We verified that the published matrices contain no missing
values, so the imputation step is a defensive no-op and cannot propagate
information across the split.


% ---------------------------------------------------------------------
% [7] ADD to end of Sec. III-E (ORA methods) — non-causal interpretation
%     (reviewer comment 4C)
% ---------------------------------------------------------------------
\textbf{Interpretation and limitations.}
The SHAP-based feature rankings and the subsequent ORA provide model-level
predictive explanations and statistical associations with annotated cell-type
signatures. These analyses do not establish causal relationships between
individual genes, molecular pathways, and tumour development. The observed
epithelial and alveolar lineage enrichment should therefore be interpreted as
functional context for the model-derived features rather than as evidence of
causal biological mechanisms. Furthermore, the significant enrichment observed
at the top-150 gene cutoff was not retained at the top-100 or top-200 cutoffs,
indicating sensitivity to the feature-selection threshold. These findings
should therefore be considered exploratory and require independent biological
validation.


% ---------------------------------------------------------------------
% [8] REPLACE the ORA result sentence (Sec. IV-E) — reviewer comment 4C
% ---------------------------------------------------------------------
% FROM:
%   ... indicates that the mRNA features selected by the multi-omics framework
%   suggests that the model-ranked mRNA features capture lineage-associated
%   molecular patterns relevant to tumor-type classification.
% TO:
The prominent representation of epithelial lineages (including urothelial,
luminal epithelial, and basal-like cells) alongside lung-specific tissues
indicates that the model-ranked mRNA features capture lineage-associated
molecular patterns relevant to tumour-type classification; this is an
association observed in the model's feature rankings and not evidence of a
causal mechanism.


% ---------------------------------------------------------------------
% [9] REPLACE abstract sentences — new CV numbers + corrected claims
%     (audit: macro-F1 superiority is NOT supported; mRNA "improvement" is NOT supported)
% ---------------------------------------------------------------------
% FROM:
%   Under a 5-fold stratified CV, XGBoost attains the best performance
%   (Accuracy: 0.9756, Macro-F1: 0.9518, Cohen's kappa: 0.9744). ...
%   recovers a significant (FDR <= 0.05) epithelial lineage signal at the
%   mRNA 150-gene cutoff. ... validating the necessity of a multi-omics approach.
% TO:
Under a 5-fold stratified cross-validation performed within the 80\% training
partition, XGBoost attains the best performance (Accuracy: 0.9725, Macro-F1:
0.9443, Cohen's $\kappa$: 0.9712). Corrected paired comparisons over 25
repeated folds show that XGBoost significantly outperforms both the linear SVM
and Logistic Regression on weighted F1 and accuracy (Nadeau--Bengio corrected
$p < 0.01$); a statistically significant macro-F1 advantage was not established.
For interpretability, the proposed work adopts an over-representation analysis
and identifies an exploratory epithelial lineage-associated enrichment signal
(FDR $\leq$ 0.05) at the top-150 mRNA feature cutoff. Leave-one-modality-out
experiments show that an mRNA-only model is significantly worse than the full
multi-omics model (macro-F1 0.9038 versus 0.9443), whereas removing any single
modality does not significantly reduce macro-F1. These results indicate that
global feature importance does not necessarily imply modality-level predictive
necessity.


% ---------------------------------------------------------------------
% [10] REPLACE conclusion — audit-corrected
% ---------------------------------------------------------------------
For the 32-class pan-cancer classification, XGBoost achieved the strongest
overall performance and the highest cross-validation stability among the
evaluated classifiers, with cross-validated weighted precision 0.9730, F1 0.9719,
accuracy 0.9725, Cohen's $\kappa$ 0.9712 and macro-F1 0.9443. Paired comparisons
over 25 repeated folds, corrected for fold dependence, showed XGBoost to be
significantly better than both the linear SVM and Logistic Regression on
weighted precision, F1, accuracy and $\kappa$ ($p \leq 0.011$); differences on
macro-F1 did not reach significance after correction. Background-constrained
over-representation analysis identified an exploratory epithelial and alveolar
lineage-associated enrichment among the top-ranked mRNA features at the
150-gene cutoff; these findings provide statistical context for the model-derived
predictions but do not establish causal biological relationships. A
leave-one-modality-out ablation showed that an mRNA-only model was significantly
worse than the full multi-omics model (macro-F1 0.9038 versus 0.9443,
$p < 0.001$), while removing any individual modality did not significantly
reduce macro-F1. The discriminative signal is therefore distributed across
modalities rather than localised in a single one, and global feature importance
should not be read as evidence of modality-level necessity.


% ---------------------------------------------------------------------
% [11] ADD to Limitations — statistical qualification (audit-corrected)
% ---------------------------------------------------------------------
After correcting for the dependence between repeated-CV folds
(Nadeau--Bengio), the advantage of XGBoost remains statistically significant on
weighted precision, weighted F1, accuracy and Cohen's $\kappa$ ($p \leq 0.011$),
but not on macro-F1 (versus SVM $p = 0.509$; versus Logistic Regression
$p = 0.137$). Claims of superiority are therefore restricted to the weighted
metrics. The fold-level variance of macro-F1 is dominated by the two smallest
classes (OV, $n = 7$ within the training partition; READ, $n = 73$), so
macro-averaged conclusions for these classes should be treated as indicative
only.

The leave-one-modality-out differences were likewise assessed with paired,
Nadeau--Bengio corrected comparisons over 25 folds. Only the mRNA-only model
differed significantly from the full model (macro-F1 0.9038 versus 0.9443,
$p < 0.001$); removing any single modality did not produce a significant change
in macro-F1, and removing CNV did not significantly change accuracy. These
ablations therefore support the claim that importance rankings do not imply
modality-level necessity, but they do not establish that any individual modality
is dispensable in general.

In addition, post-hoc ORA results are sensitive to the feature-selection cutoff:
statistically significant lineage enrichment ($\mathrm{FDR} < 0.05$) holds at
$N = 150$ SHAP-ranked mRNA features but does not persist at $N = 100$ or
$N = 200$.
